% chapters/10/section-06.tex
\section{Yang Belum Ada pada Pencarian Ini}
\label{sec:belum-ada-astar}

NPC sudah dapat mengetahui jalannya. Bagian ini menjelaskan apa yang masih
kurang agar pengetahuan itu benar-benar dipakai di dalam permainan.

\subsection{Pencarian yang Belum Terhubung dengan NPC}

Program pada bab ini mencari jalur dan memperlihatkannya, tetapi belum ada NPC
yang menempuh jalur itu. Penghubungan keduanya adalah pekerjaan Bab~XI, dan
penghubungan itu tidak sesederhana menambahkan beberapa baris.

Tiga persoalan muncul ketika pencarian jalur dipakai di dalam permainan.

\begin{enumerate}
  \item \textbf{Kapan pencarian dijalankan.} Pencarian tidak dapat dijalankan
        setiap bingkai untuk setiap NPC, sebab biayanya terlalu besar.
        Keputusan tentang kapan pencarian perlu diulang---misalnya ketika
        tujuan berpindah, atau ketika jalur yang ditempuh tidak lagi
        berlaku---merupakan inti Bab~XI.
  \item \textbf{Rintangan yang bergerak.} Peta pada bab ini tetap. Apabila ada
        benda yang berpindah menutupi jalan, jalur yang sudah dihitung dapat
        menjadi tidak berlaku di tengah jalan.
  \item \textbf{Jalur yang ditempuh benda.} Jalur berupa rangkaian petak,
        sedangkan benda bergerak dengan kecepatan dalam piksel per detik.
        Perlu aturan untuk menentukan kapan sebuah petak dianggap sudah
        dicapai.
\end{enumerate}

\subsection{Bentuk Peta yang Lebih Rumit}

Peta pada bab ini berpetak dan mengizinkan empat arah. Tiga bentuk
pengembangan berikut lazim dipakai pada permainan yang sesungguhnya.

\begin{itemize}
  \item \textbf{Delapan arah}, sehingga benda dapat bergerak diagonal. Seperti
        dijelaskan pada Bagian~\ref{sec:heuristik}, heuristiknya harus
        disesuaikan agar tidak melebihkan jarak.
  \item \textbf{Peta berbiaya berbeda}, misalnya jalan berbatu yang lebih
        lambat dilalui. Pencariannya tidak berubah, sebab A* memang sudah
        memperhitungkan biaya sisi. Yang berubah hanyalah nilai biaya pada
        fungsi penentuan tetangga.
  \item \textbf{\istilah{NavMesh}}, yaitu peta jalur berupa rangkaian bidang
        yang menutupi daerah yang dapat dilalui, bukan rangkaian petak. Cara itu
        lebih tepat untuk peta tiga dimensi, dan itulah yang dipakai Unity
        melalui \istilah{AI Navigation}, sebagaimana akan Anda lihat pada
        Bab~XII dan Bab~XIII.
\end{itemize}

\begin{catatabox}
Pemahaman tentang \istilah{NavMesh} pada dasarnya adalah
pemahaman yang sama dengan yang Anda peroleh pada bab ini: simpul, sisi, biaya,
dan pencarian. Perbedaannya terletak pada bentuk simpulnya, yaitu bidang
segitiga alih-alih petak, dan pada kenyataan bahwa pencariannya dilakukan oleh
mesin, bukan oleh kode Anda. Pengetahuan tentang cara kerja pencarian itulah
yang membuat Anda dapat menilai mengapa sebuah jalur yang dihasilkan mesin
terasa aneh, dan apa yang perlu diubah pada pengaturannya.
\end{catatabox}

\subsection{Ringkasan Perkembangan}

\begin{table}[htbp]
\centering
\caption{Bagian yang ditambahkan pada pencarian jalur pada bab berikutnya}
\label{tab:perkembangan-astar}
{\setlength{\tabcolsep}{4pt}\footnotesize%
\begin{tabularx}{\textwidth}{@{}C{1.6cm} L{5.0cm} Y@{}}
\toprule
\headrow \textbf{Bab} & \textbf{Yang ditambahkan} & \textbf{Akibatnya} \\
\midrule
Bab~XI & Penghubungan dengan mesin keadaan, penyimpanan jalur, dan pengulangan
pencarian secara berkala &
NPC mengejar dengan rute memutar, dan biaya pencariannya diatur agar tetap
ringan. \\
\addlinespace
Bab~XIII & \istilah{NavMesh} pada Unity &
Pencarian jalur diserahkan kepada mesin, sedangkan kode Anda mengatur tujuan,
luas jangkauan, dan tingginya rintangan. \\
\bottomrule
\end{tabularx}}
\end{table}

Dengan gambaran itu, tampak bahwa bab ini melengkapi lapisan navigasi yang
telah diperkenalkan pada Bab~II. Setelah bab ini, kedua lapisan sudah tersedia:
lapisan keputusan pada Bab~IX, dan lapisan navigasi pada bab ini. Bab~XI
menghubungkan keduanya, dan di sanalah NPC yang sebenarnya terbentuk.
